\chapter{仿真结果与分析}

本章基于前文建立的系统模型、SCA轨迹优化算法及目标定位算法，通过六组仿真实验系统性地评估所提UAV-ISAC系统的性能。仿真涵盖SCA算法收敛性、通信-感知权衡特性、三维轨迹优势、定位算法对比、能量预算影响以及CRB空间分布六个维度，力求从多角度验证理论分析的正确性与算法的有效性。

\section{仿真场景与参数设置}

仿真场景设定为$2000\times2000$~m\textsuperscript{2}的矩形区域，地面通信用户位于$(300,400)$~m，静止目标位于$(1843,1709,18.3)$~m。UAV从基站位置$(0,0,0)$起飞，标称飞行高度$H=150$~m，可在$[50,200]$~m范围内进行三维机动。仿真涉及的关键参数已汇总于表\ref{tab:sim_params}。

实验软件环境为Python~3.12，凸优化求解器采用MOSEK~10（兼容CLARABEL/SCS作为备选），所有实验均在配备Intel~Core~i7处理器、16~GB内存的个人计算机上完成。

\section{SCA算法收敛性分析}

本节验证第3章提出的基于逐次凸逼近（SCA）的轨迹优化算法的收敛特性。实验设置权衡系数$\eta=0.7$，每阶段SCA最大迭代次数$I_{\max}=20$，收敛容差$\varepsilon_{\text{obj}}=10^{-3}$。

\subsection{收敛行为}

图\ref{fig:s1}展示了单一阶段内目标函数、CRB及通信速率随SCA迭代次数的演化曲线。从图中可以观察到以下特征：

\begin{enumerate}
\item \textbf{目标函数单调下降}：目标函数$J(\mathbf{S})=\eta\cdot\text{CRB}-(1-\eta)\bar{R}_{\text{com}}/\kappa$在前5次迭代内快速下降，随后进入平缓收敛阶段，至第13次迭代附近满足收敛容差条件。该行为验证了SCA框架的理论收敛性质——每次迭代求解一个原问题的凸近似，并通过线搜索保证目标函数值的非增性。

\item \textbf{CRB与速率的耦合演化}：在$\eta=0.7$的权衡设置下，CRB从初始轨迹的$9.1\times10^{3}$下降至$9.3\times10^{2}$，降幅约一个数量级；同时平均通信速率从$9.8\times10^{6}$~bps提升至$1.06\times10^{7}$~bps，实现了感知与通信性能的同步改善。这一现象说明本文设计的直线初始轨迹策略（连接用户与目标估计位置的射线中点）能够为SCA迭代提供一个兼顾双目标的良好起点。

\item \textbf{收敛稳定性}：迭代过程中目标函数未出现振荡或发散现象，表明线搜索策略（候选步长集$\{1.0,0.8,\dots,0.05\}$）有效抑制了凸近似与原问题之间的偏差所可能导致的不稳定性。
\end{enumerate}

\subsection{收敛速度分析}

在$N_m=25$个航点的典型阶段中，SCA平均耗时约4.8秒（MOSEK求解器），平均迭代次数约13次。考虑到实际应用中UAV轨迹规划可以离线完成，该计算开销完全在可接受范围内。

\section{通信-感知权衡分析}

本节通过改变权衡系数$\eta\in\{0,0.5,0.8,1.0\}$，系统性地揭示通信与感知性能之间的内在权衡关系。其中$\eta=0$对应纯通信优化（$J=-\bar{R}_{\text{com}}$），$\eta=1$对应纯感知优化（$J=\text{CRB}$），中间值代表不同程度的权衡策略。

\subsection{性能指标对比}

图\ref{fig:s2_bars}以双轴柱状图的形式展示了四种策略下的最终CRB与平均通信速率。关键数据如表\ref{tab:tradeoff}所示。

\begin{table}[htbp]
\centering
\caption{不同权衡系数下的系统性能}
\label{tab:tradeoff}
\begin{tabular}{cccc}
\toprule
$\eta$ & CRB (xyz trace) & 平均速率 (bps) & 定位误差 (m) \\
\midrule
0.0（纯通信） & $1.59\times10^{5}$ & $1.32\times10^{7}$ & 541 \\
0.5（均衡）   & $9.47\times10^{2}$ & $1.08\times10^{7}$ & 155 \\
0.8（偏感知） & $9.26\times10^{2}$ & $1.06\times10^{7}$ & 100 \\
1.0（纯感知） & $9.36\times10^{2}$ & $1.05\times10^{7}$ & 274 \\
\bottomrule
\end{tabular}
\end{table}

从表\ref{tab:tradeoff}和图\ref{fig:s2_pareto}的Pareto前沿曲线中可以得出以下核心结论：

\begin{enumerate}
\item \textbf{CRB的边际递减效应}：当$\eta$从0增加至0.5时，CRB从$1.59\times10^{5}$骤降至$9.47\times10^{2}$，降幅超过两个数量级；然而当$\eta$从0.5进一步增加至1.0时，CRB仅从$9.47\times10^{2}$微降至$9.36\times10^{2}$，改善幅度不足2\%。这表明仅需在优化目标中引入中等权重的CRB项，即可获得接近最优的感知性能，该发现对实际系统设计具有重要指导意义——并非需要设定$\eta=1$才能获得良好的定位精度。

\item \textbf{通信速率的有限牺牲}：相较于纯通信策略，$\eta=0.5$方案仅牺牲约18\%的通信速率（$1.32\times10^{7}\rightarrow1.08\times10^{7}$~bps），却换取了超过两个数量级的CRB改善。这一交换的"性价比"极高，充分体现了ISAC系统中通信与感知在空间自由度上的互补特性。

\item \textbf{Pareto前沿的拐点特征}：图\ref{fig:s2_pareto}中的Pareto前沿曲线呈现明显的拐点结构——在高速率-高CRB区域（左上），曲线斜率陡峭，微小的CRB改善需付出较大的速率代价；而在低CRB区域（右下），曲线趋于平坦，进一步降低CRB所需付出的速率代价极小。该现象的内在机理在于：UAV轨迹优化的核心矛盾在于UAV需同时靠近通信用户和感知目标；当感知权重足够大时，UAV已有强烈动机接近目标，此时继续增加$\eta$不会显著改变轨迹几何，因此CRB趋于饱和。
\end{enumerate}

\subsection{轨迹几何的适应性变化}

从图\ref{fig:s2_traj}中各$\eta$值下的UAV轨迹对比可见：$\eta=0$时UAV径直飞向通信用户并在其附近盘旋，完全忽略感知需求；随着$\eta$增大，UAV轨迹逐渐向目标方向偏移；当$\eta\geq0.5$时，轨迹形态趋于稳定——UAV沿用户与目标的连线方向飞行，在两者之间取得空间折中。该轨迹形态变化与目标函数的数学结构高度一致，验证了优化框架的合理性。

\section{三维轨迹与二维轨迹对比}

本节通过对比UAV可变高度（3-D）轨迹与固定高度（2-D，$z\equiv H=150$~m）轨迹，量化三维机动对系统性能的增益。

\subsection{性能增益定量分析}

\begin{table}[htbp]
\centering
\caption{2-D与3-D轨迹的性能对比（$\eta=0.7$）}
\label{tab:2d3d}
\begin{tabular}{cccc}
\toprule
轨迹模式 & CRB (xyz trace) & 平均速率 (bps) & 定位误差 (m) \\
\midrule
2-D（固定高度） & $3.86\times10^{3}$ & $9.89\times10^{6}$ & 116 \\
3-D（可变高度） & $9.28\times10^{2}$ & $1.06\times10^{7}$ & 98 \\
\hline
3-D相对改善   & $\downarrow76.0\%$ & $\uparrow7.3\%$ & $\downarrow15.9\%$ \\
\bottomrule
\end{tabular}
\end{table}

表\ref{tab:2d3d}和图\ref{fig:s3}揭示了以下关键发现：

\begin{enumerate}
\item \textbf{高度自由度对CRB的主导性影响}：3-D轨迹的CRB为$9.28\times10^{2}$，仅为2-D轨迹CRB（$3.86\times10^{3}$）的24\%，改善幅度达76\%。这一显著增益的物理根源在于：当UAV与目标具有相似的高度时，两者的水平距离占据斜距的主导成分，FIM中z方向的偏导数项过小，导致z轴定位精度极差；3-D轨迹通过在不同高度部署悬停点，有效激发FIM各方向的响应分量，从而戏剧性地改善了CRB条件数。

\item \textbf{通信速率同步获益}：3-D轨迹的通信速率较2-D提升7.3\%。这是因为固定高度150~m限制了UAV逼近用户的能力——当UAV需同时靠近远距离目标时，较高的固定高度增大了UAV到用户的斜距。可变高度允许UAV在目标区域爬升至200~m以改善感知几何，在用户区域下降至50~m以缩短通信距离，实现双赢。

\item \textbf{定位精度的间接改善}：3-D轨迹的最终定位误差为98~m，较2-D（116~m）降低15.9\%。需要指出的是，定位误差的改善幅度小于CRB的改善幅度，这是因为定位精度还受到测量噪声、估计算法特性等多重因素的影响，但降低CRB确实为实现高精度定位提供了更有利的观测几何条件。
\end{enumerate}

从图\ref{fig:s3}的三维轨迹对比中可以直观地观察到：3-D轨迹的悬停点分布在50~200~m的高度范围内，形成立体的观测构型；而2-D轨迹的所有航点均被限制在$z=150$~m的平面上，丧失了高度维度的观测信息。这一可视化结果有力地支撑了将UAV高度纳入优化变量的必要性。

\section{MLE与EKF定位性能对比}

本节在相同的优化轨迹和测量数据下，对比最大似然估计（MLE）与扩展卡尔曼滤波（EKF）两种目标定位算法的性能，评价维度包括定位精度和计算效率。为获得统计意义的结果，实验在3个独立随机种子（噪声实现）下重复进行。

\subsection{定位精度对比}

\begin{table}[htbp]
\centering
\caption{MLE与EKF定位精度与耗时对比（均值$\pm$标准差，3次实验）}
\label{tab:mle_ekf}
\begin{tabular}{cccc}
\toprule
指标 & MLE & EKF & 比值 \\
\midrule
定位误差 (m)          & $14.9\pm4.6$ & $83.3\pm42.2$ & EKF/MLE = 5.6$\times$ \\
单阶段SCA耗时 (s)     & $5.01\pm0.18$ & $4.80\pm0.22$ & -- \\
单阶段定位耗时 (s)    & $16.94\pm3.38$ & $2.0\times10^{-4}$ & MLE/EKF $\approx 8.5\times10^{4}$ \\
\bottomrule
\end{tabular}
\end{table}

从表\ref{tab:mle_ekf}和图\ref{fig:s4}中可以得出：

\begin{enumerate}
\item \textbf{MLE的精度优势}：MLE的平均定位误差为14.9~m，仅为EKF（83.3~m）的约18\%。MLE的精度优势源于其每次定位均使用全部历史测量数据，并采用全局网格搜索策略，能够逼近负对数似然函数的全局最小值。相比之下，EKF仅利用当前时刻的单个测量进行序贯更新，信息利用效率较低。

\item \textbf{EKF的计算效率优势}：EKF单阶段定位耗时仅$2.0\times10^{-4}$~s（200~$\mu$s），约为MLE（16.9~s）的八万五千分之一。MLE的计算开销主要来自两阶段网格搜索——粗搜索覆盖整个$2000\times2000$~m\textsuperscript{2}区域（步长80~m），精搜索在粗搜索结果周围进行细化（步长2~m），总计需评估约$10^{4}$个网格点的负对数似然值。EKF仅执行常数复杂度的矩阵运算（3$\times$3矩阵求逆和乘法），计算量可忽略不计。

\item \textbf{精度-效率的工程权衡}：图\ref{fig:s4_err}展示了定位误差随测量累积的变化趋势。MLE的误差呈现单调下降，在累积足够多测量后趋于稳定；EKF的误差下降较慢且伴随较大波动，反映了其对初始估计误差和非线性测量模型的敏感性。在实际部署中，若计算资源充裕且对定位精度要求严格，MLE是更优选择；若应用场景要求实时性（如在线航迹规划），EKF以微小的精度代价换取极高的计算效率，具有显著的工程价值。
\end{enumerate}

\subsection{精度差异的机理分析}

EKF精度劣于MLE的根本原因在于非线性测量模型的线性化误差。距离测量方程$z_k=\|\mathbf{u}_k-\mathbf{t}\|+w_k$在目标位置$\mathbf{t}$处是高度非线性的。EKF在每次更新时对测量方程在预测值处进行一阶泰勒展开，忽略了高阶项。当目标估计值偏离真实值较远时（初始粗扫描后通常存在数十至数百米的偏差），线性化近似引入的截断误差显著增大，导致EKF更新方向偏离真实梯度方向。本实验中采用的Joseph协方差更新公式虽然增强了数值稳定性，但无法完全补偿线性化误差。

MLE通过在整个参数空间进行网格搜索回避了线性化问题，其负对数似然函数在全局范围内被精确评估，因此能够获得渐近有效的估计。但MLE的网格搜索策略使其实时性受限，特别是在三维搜索模式下计算复杂度为$O(N_x N_y N_z)$。

\section{能量预算影响分析}

本节通过改变总能量预算$E_{\text{tot}}\in\{15,25,35,50,65,80\}$~kJ，研究能量约束对系统综合性能的影响规律。该分析对实际UAV任务规划中的电池容量配置具有直接指导意义。

\begin{table}[htbp]
\centering
\caption{不同能量预算下的系统性能（$\eta=0.7$）}
\label{tab:energy}
\begin{tabular}{ccccc}
\toprule
$E_{\text{tot}}$ (kJ) & 完成阶段数 & CRB (xyz trace) & 平均速率 (bps) & 定位误差 (m) \\
\midrule
15  & 0 & --            & --            & $>400$ \\
25  & 1 & $1.60\times10^{3}$ & $9.94\times10^{6}$ & 286 \\
35  & 1 & $9.28\times10^{2}$ & $1.06\times10^{7}$ & 98 \\
50  & 2 & 29.5          & $8.87\times10^{6}$ & 18.5 \\
65  & 3 & 11.6          & $8.36\times10^{6}$ & 7.0 \\
80  & 4 & 10.9          & $9.27\times10^{6}$ & 5.4 \\
\bottomrule
\end{tabular}
\end{table}

表\ref{tab:energy}和图\ref{fig:s5}揭示了以下关键规律：

\begin{enumerate}
\item \textbf{能量阈值效应}：当$E_{\text{tot}}=15$~kJ时，能量在扣除粗扫描开销（约15.5~kJ）后所剩无几，无法完成任何优化阶段的轨迹飞行，系统完全失效。当能量提升至25~kJ时，系统开始运行但仅完成1个阶段，CRB为$1.60\times10^{3}$。这说明存在一个最低能量阈值（约为粗扫描能量与单阶段最低能量之和），低于该阈值的系统无法实现有效的ISAC功能。

\item \textbf{CRB的指数级改善}：能量预算从25~kJ增加至80~kJ，CRB从$1.60\times10^{3}$下降至10.9，改善幅度超过两个数量级。CRB的改善呈现前期剧烈、后期趋缓的规律——25$\rightarrow$35~kJ阶段CRB下降约42\%，35$\rightarrow$50~kJ阶段下降约97\%（从928降至29.5），而65$\rightarrow$80~kJ阶段仅下降约6\%。这一趋缓特性与多阶段优化的信息累积机制有关：后续阶段的轨迹优化基于更精确的目标估计，悬停点布局更加优化，但CRB的改善空间随阶段数增加而逐渐耗尽。

\item \textbf{定位精度的渐进收敛}：定位误差从25~kJ时的286~m单调下降至80~kJ时的5.4~m，趋近于亚米级至米级精度。值得关注的是，35$\rightarrow$50~kJ的定位误差从98~m骤降至18.5~m（降幅81\%），这与CRB在同一区间的剧烈改善高度吻合，验证了CRB作为定位精度理论下界的有效性。

\item \textbf{通信速率的非单调变化}：平均通信速率并非能量预算的单调函数——从35~kJ（$1.06\times10^{7}$~bps）到65~kJ（$8.36\times10^{6}$~bps），速率反而下降了约21\%。这一看似矛盾的现象源于多阶段累积平均的计算方式（式3-10）：后续阶段新增的航点若主要服务于感知目标（即飞向远离用户的方向以优化CRB），则平均速率可能不升反降。该现象再次揭示了通信与感知之间的内在竞争关系。
\end{enumerate}

\section{CRB空间分布分析}

本节通过在目标高度平面（$z=z_t$）上计算CRB的空间分布，直观展示轨迹优化对定位精度的空间塑造效应。

图\ref{fig:s6}以对数色阶展示了优化后轨迹对应的CRB空间分布热力图，叠加了UAV飞行轨迹和悬停点位置。

从热力图中可以观察到：

\begin{enumerate}
\item \textbf{低CRB走廊的形成}：在UAV轨迹沿线及其邻近区域形成了明显的低CRB走廊（深色区域，CRB$\approx10^{-2}\sim10^{1}$），CRB值远低于远离轨迹的地图边缘区域（CRB$\approx10^{3}$）。这条走廊的走向与UAV轨迹高度一致，从起始区域延伸至目标位置附近，表明SCA优化后的轨迹主动塑造了观测几何，使目标所在区域的定位理论精度达到最优。

\item \textbf{悬停点密度与CRB梯度的关系}：悬停点密集分布的区域（轨迹中后段，靠近目标）CRB梯度较大，呈现出悬停点周围CRB急速下降的特征；而在轨迹起点的粗扫描区域（4个规则排列的悬停点），CRB分布相对均匀但整体水平较高。这印证了CRB对观测几何敏感性的理论预期——悬停点的数量和空间分布共同决定了FIM的条件数，进而影响CRB的空间分布特征。

\item \textbf{通信用户的间接影响}：图中标注的通信用户位置（星号）并未直接落入低CRB核心区域，这反映了$\eta=0.7$权衡设置下目标函数对通信与感知的双重考量。用户位置通过影响UAV轨迹的走向（UAV需要兼顾对用户的通信覆盖），间接决定了CRB分布的整体格局。

\item \textbf{CRB的动态范围}：全图CRB值跨越约六个数量级（$10^{-2}\sim10^{4}$），这一巨大的动态范围有力说明了轨迹设计对定位精度的决定性影响——即便在相同的硬件参数和传播环境下，不同的UAV观测几何可导致定位精度的天壤之别。
\end{enumerate}

\section{本章小结}

本章通过六组仿真实验，从多个维度系统地评估了所提UAV-ISAC轨迹优化与目标定位框架的性能。主要结论归纳如下：

\begin{enumerate}
\item \textbf{算法有效性}：SCA迭代算法具有良好的收敛性和稳定性，在典型问题规模下可在数秒内完成求解，满足离线轨迹规划的实时性要求。

\item \textbf{权衡特性}：通信与感知之间存在显著但非对称的权衡关系——仅需在优化目标中赋予感知指标中等权重（$\eta\approx0.5$），即可获得接近最优的感知性能，同时保持80\%以上的通信速率。该特性对实际系统参数调优具有重要指导意义。

\item \textbf{三维机动的必要性}：将UAV高度纳入优化变量可获得76\%的CRB改善和7\%的通信速率提升，三维轨迹设计是充分发挥UAV-ISAC系统潜力的关键技术手段。

\item \textbf{定位算法选择}：MLE以较大的计算开销换取更高的定位精度，适用于离线后处理场景；EKF以微小的精度代价换取极致的计算效率，适用于在线实时应用。

\item \textbf{能量配置策略}：系统性能随能量预算增加而持续改善，但呈现边际递减规律。存在一个最小能量阈值（约25~kJ）使系统具备基本的ISAC功能，50~kJ以上的能量预算可支撑多阶段精密定位。

\item \textbf{空间特性}：优化后的UAV轨迹在其沿线塑造了显著的CRB低值区域，为目标提供有利的观测几何条件。CRB的空间分布高度非均匀，进一步凸显了轨迹优化的必要性。
\end{enumerate}

综上所述，本文提出的UAV-ISAC轨迹优化框架在理论分析（第2章至第4章）和仿真验证（本章）两个层面均展现出良好的性能，为面向6G通感一体化网络中的UAV辅助ISAC应用提供了可行的设计方案和有益的理论参考。
